\section{Introduction}

On a phone, a laptop, a car or a robot, the binding constraint for a language model is usually memory,
not arithmetic. The weights and the key--value (KV) cache have to fit in DRAM next to everything else
the device runs, and batch-1 decoding is limited by how fast those bytes can be streamed, not by FLOPs.
A model that buys capability with \emph{compute} instead of \emph{parameters} is therefore attractive at
the edge: the same weights, applied more times, cost time but no extra memory.

Looped (recurrent-depth, weight-shared) transformers do exactly this. A block of layers is applied
repeatedly, so effective depth grows while the parameter count stays fixed
\citep{dehghani2019universal,lan2020albert,giannou2023looped,saunshi2025reasoning}. Recent work scales
the idea to billions of parameters \citep{geiping2025scaling,zhu2025ouro}, lets each token choose its
own depth \citep{bae2025mixture,graves2016adaptive,banino2021pondernet}, attends across depth
\citep{yao2026arloop}, shares or replays KV caches across loops
\citep{huang2026looped,licakova2026kvreplay}, and asks which design choices make recurrence work at all
\citep{zhuang2026recurrence}. This paper asks the edge version of the question:

\begin{quote}
\emph{Can a compact recurrent language model learn to allocate its own computation and achieve better
reasoning per byte of memory than a conventional transformer?}
\end{quote}

We study it at the smallest scale where it is still meaningful (a 7.3M-parameter byte-level model,
about 14\,MB of weights in bf16) on a single 32\,GB Apple M4 laptop, and we let an autonomous swarm of
AI agents do most of the architecture search, in the style of Karpathy's
autoresearch~\citep{karpathy2026autoresearch}: a fixed harness, one editable file, a fixed five-minute
training budget, and a keep/discard rule. Our contributions:

\begin{itemize}
\item \textbf{An open, laptop-scale harness} that turns the question into two numbers per experiment:
  bits per byte on held-out web text (knowledge) and exact-match accuracy on a depth-controlled
  variable-tracking task (serial reasoning), including depths never seen in training. It enforces a
  parameter budget, reports KV-cache size, latency and peak memory, and runs any number of agents
  while serializing GPU time so the wall-clock budget stays fair (\S\ref{sec:setup}).
\item \textbf{A calibration result that changes the protocol.} At this scale, five minutes of training
  from scratch (about 250 steps) leaves reasoning exactly at chance for every model and task variant
  we tried, and we show this is slow emergence of in-context retrieval, not a bug. Experiments must
  therefore modify a pre-trained looped model rather than train new ones (\S\ref{sec:calibration}).
\item \textbf{Controlled comparisons against conventional transformers} with the same memory and with
  the same depth, trained on the same schedule for the same wall-clock time (\S\ref{sec:controls}).
\item \textbf{An agent-swarm study} of depth attention, adaptive halting, KV reuse and the training
  recipe on top of the looped model, with seed-noise calibrated keep/discard decisions and seed
  re-runs of every lane's best result (\S\ref{sec:swarm}).
\end{itemize}
